\section{Experiments}\label{sec:exp}

This section measures the implementation of \cref{sec:impl} against three schemes that prove the same statement with the same commitment machinery, and against an independent implementation.
For each measurement, every scheme runs on the same machine with the same field, NTT, Merkle trees, transcript, rate, number of queries and committed levels; they differ only in the protocol.

\subsection{Setup}\label{sec:exp:setup}

\paragraph{Machine and code.}
All runs use a cloud virtual machine with an Intel Xeon processor at $2.10$\,GHz ($2$ vCPUs, $7$\,GB RAM), the machine class of \KB{\S9}, Rust $1.97.0$, release profile with fat LTO and one codegen unit, and no \texttt{target-cpu} flag.
The libraries are \texttt{sumfri}~\sfversion{} and \texttt{kbfold}~0.3.0; the script \texttt{sumfri/rust/bench/run\_all.sh} runs every measurement of this section and writes the data of the tables to \texttt{sumfri/rust/results/}.
Runs are single-threaded.
Times are medians over $11$ runs for $n \le 16$, $5$ for $n \in \{18, 20\}$ and $3$ for $n = 22$; verification times are medians over $21$ runs.
The machine is shared; we repeated the measurements of \cref{tab:sum,tab:pq} in a second, independent run, whose figures are in the repository and lead to the same conclusions.

\paragraph{Parameters.}
The \emph{classical parameters} are those of \KB{\S9.1}: rate $\rho = 1/4$, $\ell = n - 4$ folding rounds (final table of $16$ elements), $\kappa = 148$ queries, the quadratic extension $\K$, $32$-byte salts, and committed levels $\mathcal{J} = \{0, 4, 8, \dots\}$.
By \cref{thm:sound,thm:skip} with $\delta = 3/8$, the soundness error is below $2^{-100}$ (\cref{rem:params}).
The \emph{post-quantum parameters} are those of \cref{rem:pq-params}: the same rate, $\ell$ and $\mathcal{J}$, $\kappa = 248$ queries and the quartic extension.
Tables are uniformly random over $\mathbb{F}_p$; evaluation points are uniformly random in $\F^n$.
Proof sizes are the lengths of the canonical encodings; they depend slightly on the query points, through the number of distinct opened leaves.

\paragraph{Schemes.}
\begin{itemize}
  \item \emph{Sum-FRI}: $\Pisum$, and $\Pieval$ at a random point (\cref{sec:impl}).
  \item \emph{KBFold}: the scheme of~\cite{Mkh26} (\texttt{kbfold}, kernel encoding), which proves the sum of its committed table as $2^n \tilde f(\tfrac12, \dots, \tfrac12)$ with a degree-$2$ sumcheck, and $\tilde f(z)$ at a random point.
  \item \emph{Coefficient form with sumcheck}: the coefficient-form mode of \texttt{kbfold} (\KB{\S9.1}), which commits to the Kronecker embedding of the coefficients of $\tilde f$ (after M\"obius inversion) and proves $\tilde f(z)$ with a degree-$2$ sumcheck and the classical fold, as in Reed--Solomon instantiations of BaseFold. For the sum, $z = (\tfrac12, \dots, \tfrac12)$.
  \item \emph{Quotient}: the opening of $V_f$ at $X = 1$ with a quotient and FRI (\cref{sec:impl:quotient}); it proves the sum only.
\end{itemize}
In all four schemes the commitment is one NTT of size $M$ and one salted tree of $M/2$ leaves, and the prover commits to the same number of folded words; the schemes differ in the transform before the NTT (none, kernel, M\"obius, none on a coset), in the round messages and the work that produces them, and, for the quotient, in the computation of the quotient word.
Every scheme commits to a table of $N$ field elements and proves the sum of its entries, so the four rows of \cref{tab:sum} prove the same statement about the same data.
For evaluations, Sum-FRI proves $P_f(z)$ and the other schemes $\tilde f(z)$, two different statements of the same cost class.

\subsection{The sum}\label{sec:exp:sum}

\begin{table}[t]
\centering
\small
\begin{tabular}{r r rrrr r r}
\toprule
 & Commit & \multicolumn{4}{c}{Open (ms)} & Verify & Proof \\
$n$ & (ms) & Sum-FRI & KBFold & coeff.+sc. & quotient & (ms) & (KiB) \\
\midrule
12 & 1.6 & 0.47 & 0.73 & 0.70 & 0.80 & 0.45 & 79 \\
14 & 7.2 & 2.1 & 3.2 & 3.0 & 3.1 & 0.60 & 115 \\
16 & 26.7 & 7.7 & 11.1 & 10.8 & 13.8 & 0.57 & 149 \\
18 & 136.4 & 31.9 & 43.6 & 45.3 & 61.5 & 0.90 & 193 \\
20 & 589.3 & 144.5 & 246.4 & 219.3 & 322.3 & 1.06 & 237 \\
22 & 3021.5 & 782.4 & 1150.1 & 1113.8 & 1474.4 & 1.32 & 288 \\
\bottomrule
\end{tabular}
\caption{Proving $\sum_a f(a)$ for a committed table of $N = 2^n$ elements, classical parameters, single-threaded. Commit, verify and proof size are those of Sum-FRI; the other three schemes are within the run-to-run variation of these figures, except the verifier of the quotient scheme, which is $6$--$56\%$ slower (\cref{sec:exp:sum}). The verification time at $n = 20$ is from the second run (the first gave $2.2$\,ms for Sum-FRI and $1.1$\,ms for the other schemes).}
\label{tab:sum}
\end{table}

\Cref{tab:sum} reports the four schemes on the sum.
The opening phase of Sum-FRI is $27$--$41\%$ faster than that of KBFold and $27$--$36\%$ faster than the coefficient-form scheme with sumcheck, over both runs and all sizes; it is $33$--$55\%$ faster than the quotient opening.
Commit times agree within the run-to-run variation (the four commitments perform the same NTT and the same tree; the transforms of KBFold and of the coefficient-form scheme cost $\tfrac n2 N$ additions, invisible next to the NTT), and so do proof sizes and the verification times of the two sumcheck schemes, which are dominated by the $\kappa$ authentication paths and opened cosets; the quotient verifier is $6$--$56\%$ slower, since it recomputes the quotient at the opened points.
The round messages, $\ell$ field elements for Sum-FRI and $3\ell$ for the sumcheck schemes, differ by at most $0.6$\,KiB.

The saving is the one predicted in \cref{sec:scheme:costs}.
With $z = \mathbf 1$, a restriction round of Sum-FRI costs one multiplication per entry of the table (the restriction $\mathrm{cres}_\theta$) and no multiplication for the round message (two sums), while a sumcheck round restricts two tables (the table and the equality table) and computes three products per pair.
The folds, the trees of the committed words and the openings are identical in the three schemes, and they account for the remaining time.
The quotient opening is slowest: it computes the quotient word on the whole domain, with one batched inversion of $M$ elements, before folding it.

\subsection{Evaluations}\label{sec:exp:eval}

\begin{table}[t]
\centering
\small
\begin{tabular}{r rrr r}
\toprule
 & \multicolumn{3}{c}{Open (ms)} & Proof \\
$n$ & Sum-FRI, $P_f(z)$ & KBFold, $\tilde f(z)$ & coeff.+sc., $\tilde f(z)$ & (KiB) \\
\midrule
12 & 0.75 & 0.76 & 0.75 & 79 \\
14 & 2.7 & 3.3 & 2.8 & 115 \\
16 & 12.1 & 14.1 & 13.6 & 148 \\
18 & 39.7 & 47.8 & 46.7 & 194 \\
20 & 219.3 & 256.8 & 264.0 & 236 \\
22 & 1064.7 & 1115.9 & 1123.8 & 290 \\
\bottomrule
\end{tabular}
\caption{Evaluation at a random point, classical parameters, single-threaded. Commit and verify times and proof sizes agree with \cref{tab:sum} within the run-to-run variation.}
\label{tab:eval}
\end{table}

At a random point (\cref{tab:eval}), the opening of Sum-FRI is $0$--$17\%$ faster than both sumcheck schemes.
The restriction round now needs the tensor $\tau_j$ (\cref{sec:impl:enc}): two inner products with $\tau_j$ per round, plus $N$ multiplications once to build $\tau_0$, against three products per pair and the restriction of the equality table for the sumcheck.
The gap is smaller than for the sum, and at $n = 22$ it is within the run-to-run variation.

\subsection{Post-quantum parameters}\label{sec:exp:pq}

\begin{table}[t]
\centering
\small
\begin{tabular}{r rr rr rr}
\toprule
 & \multicolumn{2}{c}{Open, sum (ms)} & \multicolumn{2}{c}{Open, evaluation (ms)} & Verify & Proof \\
$n$ & Sum-FRI & KBFold & Sum-FRI & KBFold & (ms) & (KiB) \\
\midrule
12 & 0.87 & 1.42 & 1.03 & 1.42 & 1.07 & 131 \\
14 & 3.3 & 7.1 & 5.1 & 7.6 & 1.38 & 215 \\
16 & 14.1 & 29.0 & 21.6 & 26.5 & 1.68 & 295 \\
18 & 58.2 & 116.1 & 88.1 & 121.9 & 2.00 & 390 \\
20 & 269.3 & 626.5 & 468.9 & 605.9 & 3.30 & 498 \\
22 & 1443.5 & 2469.5 & 1971.6 & 2536.5 & 2.87 & 619 \\
\bottomrule
\end{tabular}
\caption{Post-quantum parameters (quartic extension, $\kappa = 248$), single-threaded. Verify and proof size are those of Sum-FRI for the sum. Commit times are those of \cref{tab:sum} within the run-to-run variation: the commitment does not depend on the extension.}
\label{tab:pq}
\end{table}

With the post-quantum parameters (\cref{tab:pq}), every operation of the rounds is in the quartic extension, built over $\K$ by Karatsuba (three multiplications in $\K$ per product).
Sum-FRI opens the sum $34$--$57\%$ faster than KBFold over both runs, and evaluations $19$--$33\%$ faster.
Proofs are $1.7$--$2.2$ times larger than with the classical parameters, as in \KB{\S9.3}: $100$ more queries, each opening elements of $\mathbb{F}_{p^4}$.

\subsection{Where the time goes}\label{sec:exp:breakdown}

At $n = 20$ with the classical parameters, the commitment spends $251$\,ms in the NTT and $490$\,ms in the salted tree of $M/2 = 2^{21}$ leaves (medians of separate runs of each step, whose sum exceeds the median commit time of the same run, $657$\,ms, by the run-to-run variation).
The opening of the sum takes $161$\,ms and that of an evaluation $203$\,ms: the $42$\,ms difference is the cost of the tensor $\tau_0$ and of the products with it.
For all four schemes the commitment dominates the prover, and its cost is shared: Sum-FRI changes only the opening.

\subsection{An independent reference point}\label{sec:exp:whir}

\begin{table}[t]
\centering
\small
\begin{tabular}{r rrr rrr rrr}
\toprule
 & \multicolumn{3}{c}{Sum-FRI, classical} & \multicolumn{3}{c}{WHIR, unique, $k = 1$} & \multicolumn{3}{c}{WHIR, Johnson, $k = 4$} \\
$n$ & prover & verify & proof & prover & verify & proof & prover & verify & proof \\
\midrule
12 & 2.1 & 0.45 & 79 & 18.4 & 1.40 & 131 & 4.6 & 0.38 & 45 \\
14 & 9.2 & 0.60 & 115 & 87.7 & 2.30 & 199 & 19.2 & 0.53 & 56 \\
16 & 34.4 & 0.57 & 149 & 335.1 & 2.80 & 285 & 71.2 & 0.66 & 75 \\
18 & 168.3 & 0.90 & 193 & 1318 & 3.90 & 392 & 340.6 & 0.48 & 88 \\
20 & 733.8 & 1.06 & 237 & 6183 & 3.40 & 513 & 2470 & 0.54 & 108 \\
22 & 3803.9 & 1.32 & 288 & 27000 & 5.50 & 657 & 9700 & 0.54 & 121 \\
\bottomrule
\end{tabular}
\caption{Sum-FRI ($\Pisum$, prover $=$ commit $+$ open) and the WHIR implementation, same machine, single-threaded, Goldilocks with its quadratic extension, rate $1/4$, $100$ bits of security, no proof of work. Times in ms, proof sizes in KiB. WHIR prints times of one second or more with a resolution of $0.1$\,s; its figures are those of the median, by commit time, of three invocations, each with \texttt{--reps 21}. The Sum-FRI verification time at $n = 20$ is from the second run, as in \cref{tab:sum}.}
\label{tab:whir}
\end{table}

As in \KB{\S9.8}, we run the WHIR implementation~\cite{ACFY24}\footnote{Repository \url{https://github.com/WizardOfMenlo/whir}, commit \texttt{c03a4a5}, the version used in KBFold.} on the same machine with the same field, rate and security level, single-threaded and without proof of work, in two configurations: one variable per round with the unique-decoding analysis, and four variables per round with its analysis up to the Johnson bound (\cref{tab:whir}).
WHIR proves a multilinear evaluation, and the two code bases differ in many engineering choices, so \cref{tab:whir} is a reference point, not a comparison of protocols.
The Sum-FRI prover is $7$--$10$ times faster than WHIR with one variable per round, with proofs $1.7$--$2.3$ times smaller and a verifier $3$--$5$ times faster.
Against WHIR's default configuration, the Sum-FRI prover is $2$--$3.4$ times faster, while WHIR's proofs are $1.8$--$2.4$ times smaller; the verification times are comparable for $n \le 16$, and WHIR's verifier is $1.9$--$2.4$ times faster for $n \ge 18$. As in KBFold, the difference in proof size and verification time comes from the number of queries, which WHIR reduces by its analysis up to the Johnson bound.
